\documentclass[11pt]{article}
\usepackage{amsmath,amssymb,amsthm,hyperref}
\newtheorem{theorem}{Theorem}
\newtheorem{lemma}[theorem]{Lemma}
\newcommand{\cube}{\{0,1\}^n}
\begin{document}
\title{What Precision Must Supply in an Exponent Improvement}
\author{}
\date{November 23, 2026}
\maketitle

\paragraph{Status.}
This note verifies the finite-information precision claim for the
intrinsic atom chart and strengthens it to denominator information alone.
It also disproves the analogous raw-entry alphabet claim when width
is unrestricted, while distinguishing that claim from finite binary
rational precision.
It proves no improvement of unrestricted worst-case head complexity,
no superconstant saving, and no matching lower bound.
The original attention model, its head resource, and its admissibility
conditions are unchanged.

\paragraph{Precise meaning of precision.}
A coefficient alphabet is a finite set fixed before the target Boolean
function is selected. Saying that a scalar has at most $b$ bits means
that its alphabet has at most $2^b$ elements. The alphabet may depend
on $n$, the coefficient position, and the construction, but not on
the target truth table. This notion permits arbitrarily large or small
real alphabet entries. It measures information, not approximation error,
decimal digits after the point, numerical stability, or cancellation.
Without an agreed encoding, a real parameter has no intrinsic finite
bit length. The conclusions below concern the worst-case or average
available information per coefficient; they do not require every
coefficient in every representation to use that many bits.

\section{All coefficients}

\begin{theorem}[Finite encoding bound]
Let $X$ be a set of $N$ inputs, let $\Theta$ be a finite code set,
and let $S:\Theta\times X\to\mathbb R$ be any fixed decoder.
If every $f:X\to\{-1,1\}$ has a code $\theta$ satisfying
$f(x)S(\theta,x)>0$ for every $x$, then $|\Theta|\ge 2^N$.
Consequently, if the decoder has $p$ scalar coefficients, each drawn
from a fixed alphabet of size at most $2^b$, then $bp\ge N$.
More generally, coefficients using at most $b_1,\ldots,b_p$ bits
must satisfy $\sum_j b_j\ge N$.
\end{theorem}
\begin{proof}
Each code has at most one strict sign vector. Choose one representing
code for each of the $2^N$ target vectors. Two different vectors cannot
choose the same code, since at a coordinate where their signs differ
the same real number would be both strictly positive and strictly
negative. The chosen codes are therefore distinct. A product of
coefficient alphabets has cardinality at most
$\prod_j 2^{b_j}=2^{\sum_j b_j}$, which gives the result.
\end{proof}

For $H$ affine fractions and a bias, the unnormalized affine chart
has $p=1+2H(n+1)$ coefficients. A valid attention atom is a member of
this chart with additional constraints, so counting the larger chart
is a valid necessary condition. In particular, a universal construction
with a fixed $b(n)$-bit alphabet for its effective affine coefficients
satisfies
\[
  2^n\le b(n)\bigl(1+2H(n)(n+1)\bigr).
\]
If, for fixed $C>0$ and $\varepsilon>0$ and all sufficiently large $n$,
\[
  H(n)\le \frac{C2^n}{n^{1+\varepsilon}},
\]
then
\[
 b(n)\ge
 \frac{n^\varepsilon}
 {n^\varepsilon/2^n+2C(1+1/n)}
 =\left(\frac{1}{2C}-o(1)\right)n^\varepsilon.
\]
This verifies the claimed $n^\varepsilon/2$ scale when the suppressed
head-count constant is $C=1$. The constant $C$ cannot be omitted in a
literal quantitative statement.

These finite formulas use exactly $H$ coded heads. If a construction
uses a target-dependent number $h\le H$, include $h$ in the code.
There are at most $(H+1)2^{b(1+2H(n+1))}$ such codes, so the exact
inequality becomes
$2^n\le\log_2(H+1)+b(1+2H(n+1))$. In the head-count regime above,
$\log_2(H+1)=O(n)=o(2^n)$; the displayed asymptotic precision bound
is unchanged. This argument does not assume that a particular finite
alphabet contains coefficients for padding by zero heads.

The same asymptotic bound holds for alphabets on the original
\texttt{FracAtom} parameters: there are $2n+4$ parameters per atom,
namely $\eta,\delta,\gamma,\alpha$, the $n$ values of $\rho$, and the
$n$ values of $m$, plus the global bias. Use
$p=1+H(2n+4)$ in the encoding theorem. Positivity constraints only
remove codes and do not invalidate the bound. The number of matrix
entries in a redundant attention parameterization is a different
resource; the claim is about the intrinsic scalar chart or a specified
code of comparable length, not an unrestricted redundant parameter
description. In fact, \texttt{computableWithHeadsN} existentially
quantifies an unrestricted width $d$. In the literal head-local
matrix chart, the three $d\times d$ matrices and the $(n+4)$
embedding vectors give $3d^2+(n+4)d$ scalars per head, plus $d$
readout scalars and the threshold for the family. Direct finite-code
counting therefore gives $bP\ge2^n$ with
$P=H(3d^2+(n+4)d)+d+1$, not with $P=1+H(2n+4)$.
The encoding theorem applies despite the exponential in softmax,
since it places no regularity restriction on the decoder. However,
passing from a finite matrix alphabet to the atom normal form does
not preserve its alphabet size. Accordingly the unchecked claim,
if read as a width-independent lower bound on the precision of
individual matrix or embedding entries, has not been established.

For completeness, the normalized affine chart has $1+2nH$ parameters.
Write a head as $A/B$, where
$A=a_0+\sum_i a_ix_i$ and $B=d_0+\sum_i d_ix_i$ with $d_0>0$.
Set
\[
 b_i=d_i/d_0,\qquad
 \alpha=a_0/d_0,\qquad
 u_i=a_i/d_0-\alpha b_i.
\]
Then
\[
 \frac{A(x)}{B(x)}=
 \alpha+\frac{\sum_i u_ix_i}{1+\sum_i b_ix_i}.
\]
Absorb the $H$ constants $\alpha$ into the bias. This is an exact
identity; positive scaling of a denominator preserves its whole-cube
admissibility. A finite alphabet for the original chart need not remain
the same alphabet under this normalization. Each alphabet claim must
specify which chart is encoded.

\section{Denominators must carry almost the whole table}

The next result strengthens the earlier audit's exponent-two statement
to every positive exponent improvement. Numerators and the bias now
have completely unrestricted real precision.

\begin{lemma}[Sign patterns of a linear space]
If $U\subseteq\mathbb R^N$ is a linear space of dimension at most
$d$, with $1\le d\le N$, then the number of strict sign vectors
met by $U$ is at most
\[
 \sum_{j=0}^d {N\choose j}\le (eN/d)^d.
\]
\end{lemma}
\begin{proof}
The family of positive-coordinate subsets of strict vectors in $U$
cannot shatter any $d+1$ coordinates. Indeed, the image of $U$ under
restriction to those coordinates has dimension at most $d$, so some
nonzero real vector $a$ is orthogonal to that image. Prescribe the sign
of coordinate $i$ to equal the sign of $a_i$ when $a_i\ne0$, and
arbitrarily otherwise. Every vector with these strict signs has
positive inner product with $a$, a contradiction.

Here is the required finite set-family count. For a family on $m$
coordinates, split by its last coordinate. Let $A$ be the union of
the two trace families on the first $m-1$ coordinates, and let $B$
be their intersection. The original size is $|A|+|B|$. If the original
family shatters no set of size $d+1$, then $A$ shatters no such set
and $B$ shatters no set of size $d$: a set of size $d$ shattered by
$B$, together with the last coordinate, would be shattered by the
original family. Induction on $m$, with the empty-coordinate case
and the dimension-zero case both bounded by one, gives
$\sum_{j=0}^d {m\choose j}$ by Pascal's identity. Cases $d\ge m$
are bounded by $2^m$ directly.

For the final inequality put $t=d/N\in(0,1]$. Since $t^j\ge t^d$
for $0\le j\le d$,
\[
 t^d\sum_{j=0}^d {N\choose j}
 \le \sum_{j=0}^N {N\choose j}t^j
 =(1+t)^N\le e^{Nt}=e^d.
\]
Divide by $t^d$.
\end{proof}

\begin{theorem}[Information in a denominator library]
Put $N=2^n$ and $d=1+nH\le N$. Suppose that a library of $K$
ordered lists of $H$ nonvanishing affine denominators suffices to
represent every Boolean function strictly by some list, arbitrary
affine numerators, and a bias. Then
\[
 N\le \log_2 K+d\log_2(eN/d).
\]
In particular, if $H(n)\le C2^n/n^{1+\varepsilon}$ for fixed
$C,\varepsilon>0$, a universal denominator library satisfies
\[
 \log_2 K(n)\ge (1-o(1))2^n.
\]
\end{theorem}
\begin{proof}
Fix one denominator list. The evaluations of affine numerators divided
by the $h$th denominator form a linear space $U_h$ of dimension at
most $n+1$. It contains the constant functions, because the numerator
can equal the denominator. The sum of these spaces and the bias
therefore has dimension at most $1+nH=d$, including when $H=0$.
The preceding lemma bounds its strict patterns by $(eN/d)^d$.
The union over the $K$ lists has at most $K(eN/d)^d$ patterns.
Universality requires at least $2^N$ patterns; take base-two logarithms.

For the asymptotic statement, put $r=d/N$. The head bound implies
$0<r\le 2^{-n}+C/n^\varepsilon$, so $r\to0$. Since
$r\log_2(e/r)\to0$, the error term
$d\log_2(eN/d)=Nr\log_2(e/r)$ is $o(N)$. Also $d\le N$
eventually, as required for the finite estimate.
\end{proof}

If each of the $n$ normalized coefficients of each denominator has
a fixed alphabet of at most $2^{b(n)}$ values, then
$K(n)\le 2^{b(n)nH(n)}$. Hence the last theorem gives
\[
 b(n)nH(n)\ge(1-o(1))2^n,\qquad
 b(n)\ge\left(\frac1C-o(1)\right)n^\varepsilon.
\]
With unnormalized denominator coefficients there are $(n+1)H(n)$
coordinates instead, giving the same asymptotic constant. This result
allows unbounded information in all numerator coefficients. Thus a
successful exponent improvement must have the adaptive denominator
choices themselves encode almost $2^n$ bits across the worst-case
family, when measured by the size of a fixed finite library.

Allowing lists of any length at most $H$ adds at most a factor $H+1$
to this library count. The same $O(n)=o(2^n)$ logarithmic correction
therefore preserves the denominator-only asymptotic conclusion.

Admissibility has only been relaxed when counting sign spaces. Every
admissible attention denominator is nonvanishing on the entire cube,
so all of these bounds apply to the original model. No proposed
denominator is being declared admissible merely because it is positive
on a subset of inputs. Every counted output pattern is strict.

\section{Unrestricted raw width changes the information claim}

There is a stronger reason not to identify the intrinsic coefficient
precision bound with a bound on individual raw matrix entries.

\begin{theorem}[Same-head simulation by a fixed raw alphabet]
Let $A=\{0,1,-1,\sqrt2,-\sqrt2\}$. Fix $n$ and $H\ge1$.
Every Boolean function strictly computable with $H$ heads in the
repository model has a strict $H$-head representation in which every
entry of every token embedding, positional embedding, query matrix,
key matrix, value matrix, and the common readout vector belongs to
$A$, and the threshold is zero. The width is finite and unrestricted;
it may depend on the target function and its original representation.
\end{theorem}
\begin{proof}
First, the additive subgroup $D=\mathbb Z+\mathbb Z\sqrt2$ is dense
in $\mathbb R$. To see this explicitly, for any integer $M\ge1$
place the $M+1$ fractional parts of $j\sqrt2$, $0\le j\le M$,
in $M$ half-open intervals of length $1/M$. Two lie in one interval.
Their difference is a nonzero member of $D$ with absolute value less
than $1/M$, since $\sqrt2$ is irrational. Its absolute value is a
positive $a\in D$. For any $r\in\mathbb R$ an integer multiple
of $a$ lies within $a/2$ of $r$. Thus elements of $D$ approximate
any specified real arbitrarily closely. Every element of $D$ is a
finite sum of entries of $A$, by expanding the integer coefficients
as repeated copies of $1,-1,\sqrt2$, or $-\sqrt2$.

Use the exact atom normal form to write the original strict score as
$c+\sum_h\phi_h(x)$. Define
\[
 a_h=\log\alpha_h,\quad q_h=\log\gamma_h,\quad
 r_{h,i}=\log\rho_{h,i},\quad t_h=\delta_h.
\]
Select one head $h_0$ and put
\[
 v_{h,q}=\eta_h/\gamma_h+\mathbf1_{h=h_0}c,
 \qquad v_{h,i}=m_{h,i}+\mathbf1_{h=h_0}c.
\]
Since the attention weights in each head sum to one, absorbing $c$
in all position values of head $h_0$ gives the exact score
\[
 S(x)=\sum_{h=1}^H
 \frac{e^{q_h}v_{h,q}
 +\sum_{i=1}^n e^{r_{h,i}+a_hx_i}(v_{h,i}+t_hx_i)}
 {e^{q_h}+\sum_{i=1}^n e^{r_{h,i}+a_hx_i}}.
\]
This expression is continuous in all displayed real parameters for
every input, since its denominators are sums of positive exponentials.
Let $\mu>0$ be the minimum absolute score over the finite Boolean
cube. By continuity at the original parameter vector, there is an
open neighborhood on which all $2^n$ scores differ from their
original values by less than $\mu/2$. Density of $D$ supplies a
parameter vector in that neighborhood with every displayed parameter
in $D$. Denote these approximations by the same letters. Their
scores retain every strict sign, and their denominators remain
positive on the entire cube.

It remains to realize these $D$-valued logits and position values
with raw entries in $A$. For each head reserve coordinates $0$ and
$1$. For each scalar among $a_h,q_h,r_{h,i},t_h,v_{h,q},v_{h,i}$,
choose a finite sequence of elements of $A$ summing to it, and
allocate a separate block of coordinates for that sequence.
Call the blocks for $a_h,q_h,r_{h,i}$ key blocks and the remaining
ones value blocks. Choose a common finite width large enough for
the blocks of every head, padding unused coordinates by zero.

For head $h$, set the embedding of token $0$ to zero. In the token
$1$ embedding, place the sequence for $a_h$ in its key block and
the sequence for $t_h$ in its value block. In the query-token
embedding, put $1$ in coordinate $0$, the sequence for $q_h$ in
its key block, and the sequence for $v_{h,q}$ in its value block.
At input position $i$, put the sequence for $r_{h,i}$ and the
sequence for $v_{h,i}$ in their respective positional-embedding
blocks. The query positional embedding is zero. All unspecified
entries are zero. Every raw embedding entry therefore lies in $A$.

Let $W_Q$ project coordinate $0$ to coordinate $0$. Let $W_K$ sum
all key-block coordinates into coordinate $0$, with all other output
coordinates zero. Let $W_V$ sum all value-block coordinates into
coordinate $1$, with all other output coordinates zero. Every
matrix entry is either zero or one. The query vector is the unit
vector $e_0$. At position $i$, the key vector is
$(r_{h,i}+a_hx_i)e_0$ and the value vector is
$(v_{h,i}+t_hx_i)e_1$. At the query position the key vector is
$q_he_0$ and the value vector is $v_{h,q}e_1$. Thus the inner
products used by softmax are exactly the logits in the displayed
formula for $S$. The common readout $w=e_1$ and threshold zero
give exactly $S(x)$. Its signs are strict, and its denominators are
positive for every cube input. No heads were added.
\end{proof}

This theorem is a written proof, not a Lean theorem in this run.
For a fixed $n,H$, there are only finitely many target functions, so
the finitely many widths needed for all representable targets have
a finite maximum. Padding by zeros therefore even gives one common
width for that whole class. The target information is stored in the
number, placement, and signs of the raw entries, whose count is not
a resource in the definition of $H^\ast$.

The alphabet has five elements, so each raw entry has a fixed
three-bit index in this alphabet. Consequently, under the finite
alphabet information interpretation, a lower bound growing with
$n$ on the bits per raw matrix or embedding entry is false when
width is unrestricted. This does not contradict the proved bounds
for the $O(nH)$ intrinsic scalar chart: combining an unbounded
number of raw coordinates produces many possible effective logits,
denominator coefficients, and values.

There is a necessary semantic qualification. The constant $\sqrt2$
is not exactly representable by a finite binary rational. This
construction therefore does not refute a separately specified claim
that every raw entry must be a rational from a standard $b$-bit
binary format. Under that format, one must count its actual raw
parameter description and any width restriction; neither the
finite-alphabet obstruction for the intrinsic chart nor this
five-symbol simulation settles a sharper rational-precision claim.

\section{Formal status and the remaining target}

The Lean file
\texttt{formalization/HeadComplexity/Results/PrecisionBarrier.lean}
contains unconditional proofs of the following declarations:
\begin{itemize}
\item \texttt{two\_pow\_card\_le\_card\_of\_strict\_universal};
\item \texttt{two\_pow\_card\_le\_alphabet\_pow\_of\_strict\_universal};
\item \texttt{card\_le\_bits\_mul\_parameters\_of\_strict\_universal};
\item \texttt{computableWithHeadsN\_of\_finiteAlphabetStrictFracComputable};
\item \texttt{two\_pow\_le\_bits\_mul\_affine\_parameters\_of\_finiteAlphabet\_universal}.
\item \texttt{rawFracParameterScore\_point};
\item \texttt{two\_pow\_le\_bits\_mul\_raw\_parameters\_of\_finiteAlphabet\_universal}.
\end{itemize}
All names are in namespace \texttt{HeadComplexity}. The first three
prove the finite encoding argument. The fourth verifies the connection
of the auxiliary finite-alphabet predicate to the unchanged attention
model. The fifth proves the exact original-model inequality
$2^n\le b(1+2H(n+1))$ for strict representations whose effective affine
coefficients lie in a fixed alphabet of size at most $2^b$.
The last two declarations prove that the raw scalar decoder agrees
with actual atom evaluation and prove the corresponding exact bound
$2^n\le b(1+H(4+2n))$ for alphabets on the raw \texttt{FracAtom}
parameters, without assuming precision preservation through any
normalization or exponential.

The additional Lean file
\texttt{formalization/HeadComplexity/Results/FixedDenominatorLibrary.lean}
proves the following declarations:
\begin{itemize}
\item \texttt{card\_le\_finrank\_of\_linearStrictSignSets\_shatters};
\item \texttt{linearStrictSignSets\_vcDim\_le};
\item \texttt{card\_linearStrictSignSets\_le\_sum\_choose};
\item \texttt{two\_pow\_card\_le\_library\_card\_mul\_sum\_choose};
\item \texttt{two\_pow\_two\_pow\_le\_denominator\_library\_capacity};
\item \texttt{two\_pow\_two\_pow\_le\_bits\_denominator\_capacity}.
\end{itemize}
The proof restricts a linear score map to a shattered coordinate set,
applies the existing strict-orthant surjectivity theorem to bound its
cardinality by the parameter dimension, and then uses mathlib's
Sauer theorem. The final declaration is the exact finite inequality
\[
 2^{2^n}\le 2^{bnH}\sum_{j=0}^{1+nH}\binom{2^n}{j}
\]
for normalized denominator alphabets of size at most $2^b$, with
arbitrary real numerator coefficients and bias. The logarithmic
estimates and all asymptotic corollaries remain written proofs in
this note, not Lean theorems.

The same Lean file defines
\texttt{HeadComplexity.WorstCaseExponentImprovement} to be exactly
the eventual assertion
\[
 \exists\delta>0\;\exists C>0\;\exists n_0\in\mathbb N\;
 \forall n\ge n_0,\qquad
 W(n)\le C2^n/n^{1+\delta},
\]
where $W(n)$ is the original
\texttt{HeadComplexity.worstCaseHeadComplexity}. This is only a
definition of a proposition. No proof of it is supplied or claimed.
Finite precision is a necessary resource, not a decoder: the argument
does not construct an admissible mechanism that extracts the encoded
table bits. The missing mechanism remains the universal strict-sign
construction at a sublinear fraction of $2^n/n$ heads.

\end{document}
